%% ======================================================================
\section{Further topics and recent advances}
\label{chap:advances}
%% ======================================================================

\begin{objectives}
\guiding{Where does PAC-Bayes go beyond this course?}
This last section is a reading guide. It briefly presents some active research
directions --- unbounded losses and dependent data, other divergences,
information-theoretic bounds, disintegrated bounds, domain adaptation --- with
pointers to the literature. It is a good starting point for the research part of
the project.
\end{objectives}

\subsection{Beyond bounded losses and i.i.d.\ data}

All the bounds of this course assume a loss with values in $[0,1]$ and i.i.d.\
data. The boundedness enters the proofs at a single place: the control of the
exponential moment $\mathcal I_\Delta(m)$ in the third step of
Theorem~\ref{thm:general}. For unbounded losses, such as the squared loss or the
log-loss, this term must be controlled through assumptions on the tails of the
loss (sub-Gaussian, sub-exponential, or bounded moments). This is the route
followed by \citet{alquier2016properties} and \citet{germain2016pac}, whose
analysis of Bayesian inference we met in \cref{chap:posteriors}. More recently,
\citet{rodriguez2024more} and \citet{chugg2023unified} gave unified treatments,
the latter based on nonnegative supermartingales, which also yield
\emph{anytime-valid} bounds, holding simultaneously for all sample sizes. This property matters in online learning and
sequential decision-making, where the number of observations is not fixed in
advance. Martingale techniques similarly extend PAC-Bayes to dependent data and
to reinforcement learning \citep{seldin2012pac,chugg2023unified}.

\subsection{Other divergences and tighter inequalities}

The KL divergence is not an intrinsic feature of PAC-Bayes: it comes from the
change of measure \eqref{eq:dv}. Other changes of measure, based for instance
on Hölder's inequality, yield bounds involving other divergences, such as Rényi
divergences \citep{begin2016pac}, which offer different trade-offs between the
complexity term and the confidence term. On the concentration side, the kl inequality is optimal
for Bernoulli variables but not for losses with a small variance, and
Bernstein-type inequalities exploit this variance
\citep{tolstikhin2013pac,mhammedi2019pac}. For majority votes, the relevant
excess losses take only three values, and the split-kl inequality of
\citet{wu2022split} exploits this structure. On the prior side, data-dependent
priors remain one of the most effective levers
\citep{parrado2012pac,dziugaite2018data,perez2021tighter}, as we saw in
\cref{fig:ddprior}.

\subsection{Information-theoretic generalization bounds}

PAC-Bayesian bounds control, with high probability, the average risk under a
posterior that depends on the sample. A closely related family of results
bounds the \emph{expected} generalization gap of a randomized algorithm by the
\emph{mutual information} $I(S;h)$ between the sample and the output of the
algorithm \citep{xu2017information}. The connection is direct: averaging a
PAC-Bayesian bound over $S$ and choosing as prior the average posterior
$P=\E_SQ_S$, which is the best data-independent prior in expectation, turns
$\E_S\KL(Q_S\|P)$ into $I(S;h)$. The monograph of \citet{hellstrom2025generalization}
presents both points of view in a unified way.

\subsection{Deterministic predictors: disintegration}

A recurring limitation of this course is that PAC-Bayes certifies a
\emph{randomized} predictor, and we had to work to transfer the guarantee to the
majority vote. Disintegrated bounds \citep{blanchard2007occam,rivasplata2020pac}
take another route: they hold with high probability over the joint draw of the
sample and of a \emph{single} hypothesis $h\sim Q_S$, so that the one model that
is actually deployed is certified. \citet{viallard2024general} give a general
framework and practical algorithms for such bounds. Applied to majority votes whose
weights are drawn once from a distribution, as in the stochastic votes of
\cref{chap:algos}, this approach certifies the vote that is actually used.

\subsection{Domain adaptation and transfer}

When the training distribution (the source) differs from the test distribution
(the target), the risk of a majority vote on the target can be related to its
risk on the source and to a \emph{divergence between the two domains}. PAC-Bayes
is particularly well suited to this problem because the disagreement $\dD(Q)$
can be computed on \emph{unlabeled} target data.
\citet{germain2013pac,germain2020pac} define a domain divergence as the difference
between the disagreements of the voters on the source and on the target, and
derive PAC-Bayesian domain adaptation bounds together with algorithms that
minimize them. Similar ideas are used for multi-view learning and for fairness
constraints.

\subsection{Open questions}

Several questions raised in this course remain open. The first concerns the
relation between tightness and learning. The tightest certificates, such as the
kl bounds or the CCPBB bound, are not always the best learning objectives: we
saw that minimizing the first-order bound gives an excellent bound on the Gibbs
risk and a poor vote. Understanding which bound should be minimized to obtain
the best \emph{predictor}, and not only the best certificate, is largely open.
The second concerns scale. How should priors be designed for boosted ensembles
with thousands of trees, or for ensembles of large neural networks, so that the
bounds remain non-vacuous? Compression-based priors \citep{lotfi2022pac} give a
partial answer. The third concerns diversity. Second-order bounds formalize the
role of diversity in a vote, and one may hope to design ensemble algorithms that
optimize the strength--diversity trade-off explicitly, with guarantees, beyond
the implicit randomization of random forests.

\begin{exercises}
\begin{exercise}
\textbf{(Comprehension)} A practitioner collects examples one at a time, recomputes
Seeger's bound with $\delta=0.05$ after each new example, and stops as soon as the
certificate is below a target. Explain why the reported certificate is not guaranteed
to hold with probability $0.95$. Show that applying the bound for sample size $m$ with
confidence $\delta_m=\frac{\delta}{m(m+1)}$ gives a bound valid simultaneously for all
$m\geq8$, and compute the extra cost in nats and the resulting Seeger bound for
$m=1000$, $\RS(\GQ)=0.1$, $\KL(Q\|P)=5$, $\delta=0.05$.
\end{exercise}
\begin{exercise}
\textbf{(Proof)} Let $\Hcal$ be finite, $S\mapsto Q_S$ a learning rule, and write
$\bar Q=\E_S Q_S$ for the average posterior. Show that, for any prior $P$,
$\E_S\KL(Q_S\|P)=\E_S\KL(Q_S\|\bar Q)+\KL(\bar Q\|P)$. Deduce that $\bar Q$ minimizes
$\E_S\KL(Q_S\|P)$ over data-independent priors, and that the minimum equals the mutual
information $I(S;h)$ between the sample and a hypothesis $h\sim Q_S$.
\end{exercise}
\begin{exercise}
\textbf{(Comprehension)} Explain why the disagreement $\dD(Q)$ of a vote can be
estimated on unlabeled target examples, whereas its Gibbs risk and its joint error cannot.
Suppose that the disagreement measured on unlabeled target data is $0.30$ and that the
target joint error is (somehow) known to be $0.05$: compute the target Gibbs risk and the
oracle C-bound on the target risk of the vote.
\end{exercise}
\begin{exercise}
\textbf{(Comprehension)} A certificate $\mathcal B$ holds for the Gibbs risk $\RD(\GQ)$
with probability at least $1-\delta$. A single network $h\sim Q$ is drawn once and deployed.
Is $\RD(h)\leq\mathcal B$ guaranteed with high probability? Using Markov's inequality over
$h\sim Q$, show that $\RD(h)\leq\mathcal B/\eta$ with probability at least $1-\delta-\eta$, and
compute this guarantee for $\mathcal B=0.05$ and $\eta=0.1$. What do disintegrated bounds
improve?
\end{exercise}
\end{exercises}

\subsection*{Further reading}
\addcontentsline{toc}{section}{Further reading}

For a first contact with PAC-Bayes, the tutorial of \citet{langford2005tutorial}
remains an excellent entry point, and the surveys of \citet{guedj2019primer} and
\citet{alquier2024user} cover the more recent developments, the latter from a
statistical perspective. For majority votes, the main references are
\citet{germain2015risk} for the C-bound, \citet{masegosa2020second} for the tandem
bound, \citet{wu2021chebyshev} for the Chebyshev--Cantelli family and
\citet{viallard2021self} for self-bounding algorithms. The monograph of
\citet{catoni2007pac} and the recent monograph of \citet{hellstrom2025generalization}
are more advanced.
